\begin{solapartat}
Només s'indueix una \textit{fem} sobre l'espira quan el flux del camp magnètic que travessa l'espira varia amb el temps, \punts{0,1} per tant es començarà a produir una \textit{fem} quan el punt A comenci a endinsar-se en la regió on hi ha el camp magnètic \punts{0,1} i això es produirà a partir de: $\frac{6\ \mathrm{m}}{2\ \mathrm{m/s}} = 3\ \mathrm{s}$ \punts{0,1}. A partir d'aquest instant el costat horitzontal del triangle s'endinsa com:
\[ d(t) = v(t - 3) \quad \text{\punts{0,1}} \]
Al ser una triangle rectangle isócels l'àrea que s'endinsa dintre del camp es:
\[ A(t) = \frac{1}{2}\,[v(t - 3)]^2 \quad \text{\punts{0,1}} \]
El flux de camp magnètic serà:
\begin{align*}
\Phi(t) &= \frac{1}{2}\,[v(t - 3)]^2\,B \quad \text{\punts{0,1}}\\
\varepsilon &= -\frac{\mathrm{d}\Phi}{\mathrm{d}t} = -v^2\,(t - 3)\,B \quad \text{\punts{0,2}}
\end{align*}
El enunciat ens diu que per t= 4s tenim:
\[ 160\ \mathrm{V} = |-2^2\,(4 - 3)\,B\,| \Rightarrow B = 40\ \mathrm{T} \quad \text{\punts{0,2}} \]
\end{solapartat}

\begin{solapartat}
Entre $t = 0$ i $t = 3\mathrm{s}$, l'espira es troba integrament fora del abast del camp magnètic, per tant la \textit{fem} induïda serà nul·la. \punts{0,1} Entre $t = 3\mathrm{s}$ i $t = 5\mathrm{s}$ la \textit{fem} augmenta linealment fins arribar al seu valor màxim. \punts{0,1} Seguint el conveni de la regla de la ma dreta el sentit del corrent en aquesta zona serà antihorari. \punts{0,2} A partir d'aquest instant el flux del camp és constant i per tant no es genera cap \textit{fem}. \punts{0,2} La gràfica per tant serà:
\figura[0.5]{sol_fig1.png}
\punts{0,4}
\end{solapartat}
